\documentclass{article}
\usepackage[sfdefault]{FiraSans}
\usepackage[T1]{fontenc}
\usepackage[utf8]{inputenc}
\usepackage{algorithm}
\usepackage{algorithmic}
\usepackage{amsfonts}
\usepackage{amsmath}
\usepackage{amssymb}
\usepackage{booktabs}
\usepackage{enumitem}
\usepackage{float}
\usepackage{geometry}
\usepackage{graphicx}
\usepackage{hyperref}
\usepackage{lipsum}
\usepackage{longtable}
\usepackage{microtype}
\usepackage{multirow}
\usepackage{natbib}
\usepackage{pdfpages}
\usepackage{pgfgantt}
\usepackage{setspace}
\usepackage{subcaption}
\usepackage{tabularx}
\usepackage{tikz}
\usepackage{xcolor}

% color scheme
\definecolor{primary}{RGB}{0,120,215}
\definecolor{secondary}{RGB}{255,87,34}
\definecolor{background}{RGB}{245,245,245}

% page layout
\geometry{ letterpaper, left=1in, right=1in, top=1in, bottom=1in, }

% hyperlink colors
\hypersetup{ colorlinks=true, linkcolor=primary, filecolor=secondary,
    urlcolor=primary, }

% page color
\pagecolor{background}

% no paragraph indentation
\setlength{\parindent}{0pt}

\begin{document}

\setstretch{1.15}
\definecolor{blue}{RGB}{0,62,126}

\newcommand{\mytitlepage}[2]{
    \thispagestyle{empty}

    \begin{tikzpicture}[remember picture, overlay]
        \node [inner sep=0pt] at (current page.center) {#1}; { \node [
            anchor=center, inner sep=1.25cm, rectangle, fill=blue!70!white, fill opacity=0,
            text opacity=1, minimum height=0.2\paperheight, minimum width=\paperwidth, text
            width=0.8\paperwidth, font=\fontfamily{pnc}\selectfont ] at (current
        page.center) {#2}; } \node [anchor=south east, outer sep=3pt] at (current
        page.south east) {\includegraphics[width=0.33\paperwidth]{images/logo.png}};
    \end{tikzpicture}

    \newpage
}

{ \mytitlepage{\includegraphics[width=\paperwidth]{images/background.png}} {
        \centering\fontfamily{phv}\selectfont {\Huge \bfseries Real-Time Brand
            Sentiment Analysis on Social Media Using Deep Learning \par}
        \vspace{8pt}
        {\begin{center}
                \begin{tabular*}{\textwidth}{@{\extracolsep{\fill}}c c c}
                    {\LARGE Jonathan Agustin} & {\LARGE Fernando Calderon} & {\LARGE Juliet Lawton}
                \end{tabular*}
            \end{center}}
        \vspace{24pt}
        {\LARGE\bfseries \par} } }

\pagenumbering{arabic}

\section{Introduction}

Businesses increasingly leverage social media platforms to track public
perception of their brands. Many use Text Classification, Sentiment Analysis,
and Natural Language Processing (NLP) to extract subjective information about
specific brands from the vast corpus of unstructured social media text data. We
propose to construct a system that includes training, evaluating, and deploying
Machine Learning (ML) models. The model will be trained to pinpoint a brand
within a text and categorize sentiment towards the brand as positive, negative,
or neutral. The goal is to empower businesses with real-time insights into
public sentiment about their brands, enabling businesses to quickly react and
adapt to shifts in public sentiment. Ultimately, the proposed system aims to
enhance the ability of businesses to maintain a positive brand image and gain a
competitive edge.

\section{Problem Statement}

Consider a set $T = \{t_1, t_2, \ldots, t_n\}$, where each element $t_i$
represents a unique social media post. Define $B = \{b_1, b_2, \ldots, b_m\}$
as a finite set that symbolizes distinct brands, with $b_0$ as a special
element signifying the absence of a recognizable brand. Define $S = \{-1, 0,
    1\}$ as a finite set that represents sentiment labels, with $s_0$ as a special
element that denotes the absence of a discernible sentiment. The sentiment
labels are defined as follows:
\begin{align*}
    -1 & = \text{negative sentiment} \\
    0  & = \text{neutral sentiment}  \\
    1  & = \text{positive sentiment}
\end{align*}

The first function $f_1: T \rightarrow B \cup \{b_0\}$ is defined as:

\[
    f_1(t_i) =
    \begin{cases}
        b_j & \text{if } t_i \text{ mentions brand } b_j \\
        b_0 & \text{otherwise}
    \end{cases}
\]

The second function $f_2: B \cup \{b_0\} \times T \rightarrow S \cup \{s_0\}$
is defined as:

\[
    f_2(b_j, t_i) =
    \begin{cases}
        s_k & \text{if } t_i \text{ expresses sentiment } s_k \text{ towards
        brand } b_j                                                          \\
        s_0 & \text{otherwise}
    \end{cases}
\]

The overall task is a composition of functions:

\[
    f = f_2 \circ (f_1 \times id_T)
\]

where $\circ$ denotes function composition and $\times$ denotes the Cartesian
product. This composition maps each post to a brand (or $b_0$) and then assigns
a sentiment label to each brand-post pair. The identity function $id_T$
maintains the original social media post in the pair, ensuring the sentiment
label is assigned to the correct post-brand pair. The output of this
composition is a sentiment label for each social media post.\\
\\
\textbf{Limitations.} The model's primary limitation is its unimodal nature. It
processes only text data, ignoring non-textual elements like images, videos,
and audio clips. Often found in social media posts, these elements can contain
significant sentiment-related information. This information can support,
contradict, or add nuance to the sentiment expressed in the text. Audio clips,
for example, can provide sentiment-related cues through tone, pitch, and other
features. The proposed system, however, will identify this multimodal
information and discard the input completely. This removal may lead to the loss
of valuable information and less accurate sentiment analysis. Future work
should consider expanding the model to a multimodal framework to provide a more
comprehensive sentiment analysis of social media posts.

\section{Methodology}

The work is divided into four stages: assessment, development, deployment, and
evaluation. Each stage informs and refines previous stages, creating a feedback
loop that continuously improves the overall system.

\begin{itemize}[leftmargin=*]

    \item \textbf{Assessment.} The stage involves selecting and testing
          appropriate models and datasets. The specific requirements of the sentiment
          analysis task, the characteristics of the available data, and the nature of the
          problem guide this selection process. The chosen models must effectively handle
          the unstructured and noisy data typical of social media and accurately capture
          and predict sentiments. The selected datasets must accurately represent the
          social media environment and provide a diverse view of expressed sentiments.

    \item \textbf{Development.} This stage preprocesses data in real-time to
          simulate a live data stream and training the selected models. The unstructured
          and noisy nature of social media data necessitates the transformation of text
          into a structured format suitable for modeling. This process involves handling
          invalid values and discarding non-textual content such as images, videos, or
          audio. Automation ensures efficiency in this process. The selected models are
          then adapted to the preprocessed dataset and trained to identify a brand or
          align with the sentiment expressed in the social media text.

    \item \textbf{Deployment.} This stage deploys the models to process social
          media data in near real-time. This stage requires building infrastructure to
          support the models and integrating the models with social media platforms such
          as Twitter where data will be streamed.

    \item \textbf{Evaluation.} This stage measures the models' ability to
          identify brands and predict sentiment in social media. For a time period,
          social media posts are collected and manually labeled for brand and sentiment.
          Standard classification metrics such as Accuracy, Precision, Recall, and F1
          Score are used to evaluate the models. The models are then ranked by
          performance.

\end{itemize}

\section{Expected Behaviors \& Outcomes}

We expect our models to demonstrate a nuanced understanding of human language,
accurately identifying brands and sentiments. Despite the inherent challenges
associated with sentiment analysis and brand identification in text data, we
anticipate that our models will achieve an accuracy rate exceeding 50\%. We
foresee our models to have real-world applicability, serving as a valuable tool
for businesses to accurately analyze sentiment from social media data. The
proposed system is expected to provide actionable insights that enable
companies to make informed decisions.

\newpage

\begin{thebibliography}{}

    \bibitem[Devlin et al., 2018]{devlin2018bert}
    Devlin, J., Chang, M. W., Lee, K., \& Toutanova, K. (2018). \newblock BERT:
    Pre-training of Deep Bidirectional Transformers for Language Understanding.
    \newblock {\em arXiv preprint arXiv:1810.04805}.

    \bibitem[Liu et al., 2019]{liu2019roberta}
    Liu, Y., Ott, M., Goyal, N., Du, J., Joshi, M., Chen, D., \ldots \&
    Stoyanov, V. (2019). \newblock RoBERTa: A Robustly Optimized BERT Pretraining
    Approach. \newblock {\em arXiv preprint arXiv:1907.11692}.

    \bibitem[Radford et al., 2019]{radford2019language}
    Radford, A., Wu, J., Child, R., Luan, D., Amodei, D., \& Sutskever, I.
    (2019). \newblock Language Models are Unsupervised Multitask Learners.
    \newblock {\em OpenAI Blog, 1(8)}.

    \bibitem[Sanh et al., 2019]{sanh2019distilbert}
    Sanh, V., Debut, L., Chaumond, J., \& Wolf, T. (2019). \newblock
    DistilBERT, a distilled version of BERT: smaller, faster, cheaper and lighter.
    \newblock {\em arXiv preprint arXiv:1910.01108}.

    \bibitem[Hutto \& Gilbert, 2014]{hutto2014vader}
    Hutto, C.J., \& Gilbert, E. (2014). \newblock VADER: A Parsimonious
    Rule-based Model for Sentiment Analysis of Social Media Text. \newblock {\em
        Eighth International Conference on Weblogs and Social Media (ICWSM-14). Ann
        Arbor, MI, June 2014}.

\end{thebibliography}

\thispagestyle{empty}

\end{document}
